%=====================================================================
\chapter{Qualitative Data Analysis in Practice}\label{ch:qda}
%=====================================================================
\locator{4}{Part~4 --- Measurement, Data and Analysis}

\begin{outcomes}
By the end of this chapter you will be able to:
\begin{tight}
  \item \textbf{Organise} a qualitative analysis so it can be audited.
  \item \textbf{Use} open-source coding tools without letting the tool shape
    the analysis.
  \item \textbf{Establish} trustworthiness using criteria appropriate to
    qualitative work.
  \item \textbf{Map} your practices onto the criteria your dissertation is
    formally assessed against.
  \item \textbf{Write} up qualitative findings so that evidence supports
    interpretation visibly.
\end{tight}
\end{outcomes}

\begin{prereq}
\chref{qualitative} for coding, saturation and the reliability debate. This
chapter is the practical and evaluative half: tools, audit trail, and the
standard your work will be judged by.
\end{prereq}

\begin{keyterms}
audit trail $\bullet$ trustworthiness $\bullet$ credibility $\bullet$
transferability $\bullet$ dependability $\bullet$ confirmability $\bullet$
thick description $\bullet$ negative case analysis $\bullet$ peer debriefing
$\bullet$ member checking $\bullet$ reflexivity $\bullet$ CAQDAS
\end{keyterms}

\section{Rigour Without Replication}

Quantitative rigour rests on replication and quantified uncertainty. Neither is
available here: another researcher analysing your transcripts would not produce
identical themes, and under an interpretivist position they are not expected to.

That does not leave qualitative work without standards. It relocates them.
Rigour becomes \textbf{transparency}: a reader must be able to see how you got
from data to claim, and to judge whether the route was reasonable.

\begin{center}
\small
\begin{longtable}{@{}L{3.0cm}L{4.4cm}L{6.6cm}@{}}
\toprule
\textbf{Criterion} & \textbf{Quantitative analogue} & \textbf{Established by} \\
\midrule
\endfirsthead
\toprule
\textbf{Criterion} & \textbf{Analogue} & \textbf{Established by} \\
\midrule
\endhead
\bottomrule
\endfoot
Credibility
  & Internal validity
  & Prolonged engagement, triangulation, member checking, negative case
    analysis, peer debriefing \\
Transferability
  & External validity
  & Thick description of context, so a reader can judge applicability to
    theirs. \emph{You} do not claim transfer; you enable the reader to assess it \\
Dependability
  & Reliability
  & Audit trail: decisions, dated, with reasons, so the process can be followed \\
Confirmability
  & Objectivity
  & Findings traceable to data rather than to the researcher's preferences;
    reflexivity about position \\
\end{longtable}
\end{center}

\begin{conceptbox}[These map onto the regulations almost line by line]
\reg{PhD Regs 2024}{\S14.3(A)} is, in substance, a trustworthiness checklist,
and reading it that way tells you what your examiner is looking for:

\begin{center}
\small
\begin{tabular}{@{}L{6.6cm}L{7.4cm}@{}}
\toprule
\textbf{The clause asks} & \textbf{Which is} \\
\midrule
``Does the research illuminate the subjective meaning, actions and contexts of
those being researched?'' & Credibility \\
``Is there evidence of adaption and responsiveness of the research design to
the circumstances of real-life social settings?'' & Dependability, and the
theoretical sampling of \chref{qualitative} \\
``Is the description detailed enough to allow the researcher or reader to
interpret the meaning and context?'' & Transferability, via thick description \\
``How are the different sources of knowledge about the same issue compared and
contrasted?'' & Triangulation \\
``How does the research move from a description of the data through quotation
or examples to an analysis and interpretation of the meaning and significance
of it?'' & Confirmability, and the analytic ladder below \\
\bottomrule
\end{tabular}
\end{center}
\end{conceptbox}

\section{The Audit Trail}

\begin{center}
\small
\begin{tabular}{@{}L{3.4cm}L{10.9cm}@{}}
\toprule
Raw data          & Recordings, transcripts, documents, field notes ---
                    preserved unaltered and separately from analysis \\
Data reduction    & Coded transcripts, each code applied and traceable \\
Data reconstruction & Category structures, theme definitions, and the versions
                    they went through \\
Process notes     & Dated memos: decisions, reasons, changes of mind \\
Materials         & Interview guide, consent form, recruitment text, protocol
                    revisions \\
Reflexivity notes & Your position, prior beliefs, and where you noticed them
                    influencing interpretation \\
\bottomrule
\end{tabular}
\end{center}

\begin{alertbox}[Version control your qualitative analysis]
Keep the trail in a Git repository, exactly as you would code. Commit after
each coding session with a message saying what changed and why.

\medskip
This gives you a dated, immutable record of when every code was created and
every theme revised --- which is precisely what dependability requires, and
precisely what a supervisor or examiner needs to see. Keep the repository
private and be careful about what goes into it: raw transcripts are personal
data (\chref{ethics}), so commit de-identified transcripts and hold the key
separately.
\end{alertbox}

\section{Tools}

\begin{center}
\small
\begin{tabular}{@{}L{2.6cm}L{3.0cm}L{8.5cm}@{}}
\toprule
\textbf{Tool} & \textbf{Licence} & \textbf{Notes} \\
\midrule
Taguette   & Open source
           & Highlight-and-tag, browser-based, exports to CSV and HTML.
             Sufficient for most thesis-scale work, and free \\
QualCoder  & Open source
           & Desktop; handles text, image, audio and video; supports coder
             comparison and simple matrices \\
NVivo, ATLAS.ti, MAXQDA & Proprietary
           & Powerful, expensive. Check what your institution licenses before
             paying \\
Spreadsheet & ---
           & Entirely adequate for one or two transcripts. Do not let tool
             selection delay analysis \\
\bottomrule
\end{tabular}
\end{center}

\begin{pitfall}[The tool is not the method]
Software stores codes, retrieves segments and counts occurrences. It does not
interpret, and three specific distortions follow from forgetting that:

\begin{tight}
  \item \textbf{Code proliferation.} Creating a code is one click, so analyses
    end up with 300 codes and no structure. Codes should be consolidated
    deliberately, and the consolidation is analysis.
  \item \textbf{Counting as evidence.} ``This code appeared 47 times'' invites
    a quantitative reading of purposively sampled data. Frequency describes
    your corpus; it estimates nothing (\chref{mixed}).
  \item \textbf{Fragmentation.} Retrieving all segments for a code strips them
    from the narrative they belonged to. Re-read whole transcripts periodically,
    or you will build themes out of decontextualised fragments.
\end{tight}
\end{pitfall}

\section{From Code to Claim}

The move that \S14.3(A) asks about --- ``from a description of the data through
quotation to an analysis and interpretation of the meaning and significance''
--- is the hardest thing in qualitative writing. It has four rungs, and weak
theses stop at the second.

\begin{center}
\small
\begin{tabular}{@{}L{1.0cm}L{3.4cm}L{9.9cm}@{}}
\toprule
1 & Data          & ``I just approve it if the tests pass, honestly.'' \\
2 & Code          & \emph{Deferring to automated checks} \\
3 & Theme         & \emph{Automation as licence to disengage} --- reviewers
                    treat passing CI as sufficient evidence of correctness,
                    substituting the tool's judgement for their own \\
4 & Interpretation & Where automated checks are visible in the review
                    interface, they reframe review from evaluation to
                    confirmation. This suggests the placement of automated
                    signals is a design decision with consequences for review
                    depth --- testable, and actionable \\
\bottomrule
\end{tabular}
\end{center}

\begin{conceptbox}[Using quotations as evidence rather than decoration]
\begin{tight}
  \item \textbf{Quote to support a claim you have just made}, not to make it.
    The analytic sentence comes first; the quotation demonstrates it.
  \item \textbf{Attribute} --- P7, senior developer, Firm B --- so a reader can
    see whether one participant is carrying the whole theme.
  \item \textbf{Report distribution honestly.} ``Three of fourteen
    participants'' where that is the case. A theme resting on one vivid
    quotation should say so.
  \item \textbf{Include disconfirming quotations.} A theme presented with only
    supporting evidence looks selected, because it may have been. Showing where
    it did not hold, and saying why, is what negative case analysis produces.
\end{tight}
\end{conceptbox}

\section{Techniques That Strengthen Credibility}

\begin{center}
\small
\begin{tabular}{@{}L{3.4cm}L{10.9cm}@{}}
\toprule
Negative case analysis
  & Actively hunt for data contradicting your theme. Either the theme is
    revised, or its scope condition is stated. This is falsification
    (\chref{philosophy}) in qualitative form, and it is the most persuasive
    single practice available \\
Member checking
  & Return findings to participants. Their disagreement is data --- but it does
    not automatically override your analysis, since participants may not share
    your analytic frame. Report what they said and what you did about it \\
Peer debriefing
  & A colleague outside the study challenges your interpretations. Record the
    challenges and your responses in the audit trail \\
Prolonged engagement
  & Enough time in the setting to distinguish the typical from the incidental \\
Reflexivity
  & State your position: prior experience, relationship to participants,
    expectations. Not confession --- a description of the instrument, which in
    this method is you \\
\bottomrule
\end{tabular}
\end{center}

\begin{traditions}{qualitative analysis}
\tradEMP{Sees it as generating hypotheses and defining constructs for later
measurement. A partial reading, but a real contribution.}
\tradDSR{Uses it to characterise the problem before designing and to
understand appropriation afterwards --- especially uses the designers did not
anticipate.}
\tradQUAL{Its home. Trustworthiness replaces validity and reliability, and
transparency replaces replication.}
\tradFORM{Sceptical by disposition, but qualitative studies of how
practitioners reason about formal notation have explained non-adoption better
than any formal argument has.}
\end{traditions}

\begin{lab}[Build an auditable analysis]
Using the transcript you coded in \chref{qualitative}:

\begin{tightnum}
  \item Set up a Git repository with the six audit-trail components as
    directories. De-identify before committing anything.
  \item Import into Taguette or QualCoder and apply your codebook.
  \item Consolidate to at most 25 codes. Write one memo explaining each merge.
  \item Build one theme up all four rungs of the ladder, ending with an
    interpretation that predicts something.
  \item Conduct negative case analysis on that theme: find contradicting data,
    or state the scope condition it required.
  \item Write the trustworthiness paragraph, mapping each practice onto the
    matching clause of \S14.3(A).
\end{tightnum}
\end{lab}

\begin{checkpoint}
\begin{tightnum}
  \item Why does qualitative rigour rest on transparency rather than
    replication?
  \item ``This code appeared 47 times'' --- what is wrong with offering this as
    evidence, and what would be legitimate?
  \item Give an example from your area of moving from theme to interpretation,
    where the interpretation predicts something.
  \item A participant disagrees with your central finding during member
    checking. Give a case where you should revise and one where you should not.
\end{tightnum}
\end{checkpoint}

\begin{chaptersummary}
\begin{tight}
  \item Rigour here is transparency: credibility, transferability,
    dependability, confirmability --- which \S14.3(A) asks for almost clause by
    clause.
  \item Keep an audit trail in version control, de-identified, committed after
    each session.
  \item Open-source tools are sufficient. The tool stores and retrieves; it
    does not interpret, and it encourages proliferation, counting and
    fragmentation.
  \item Climb all four rungs: data, code, theme, interpretation. Stopping at
    themes is the commonest weakness.
  \item Quote to support a claim, attribute, report distribution honestly, and
    include disconfirming evidence.
  \item Negative case analysis is the single most persuasive practice
    available.
\end{tight}
\end{chaptersummary}

\begin{reviewq}
  \item Trustworthiness criteria were proposed as parallels to validity and
    reliability. Does the parallel help qualitative research gain acceptance,
    or does it concede the quantitative framing? Argue a position.
  \item Version-controlled audit trails would make qualitative analysis far
    more transparent, yet almost nobody does it. Identify the practical and
    epistemic obstacles, and say which is the real one.
  \item Member checking assumes participants are well placed to validate
    analysis. When is that assumption false, and what should replace it there?
  \item Large language models are now used to code qualitative data at scale.
    Given \reg{PhD Regs 2024}{\S32(IV)}, which forbids AI from analysing and
    interpreting data, what --- precisely --- could be automated without
    breaching it, and would the result still be qualitative analysis?
    (\chref{genai} continues this.)
\end{reviewq}
